\section{Monadic types, sharp bounds and graph reachability}\label{sec:expressive}
\subsection{Counting colours in a unary signature}
Consider $p$ unary predicates, $z$ nullary predicates, equality, and no constants or positive-arity functions. A colour $c<2^p$ is the $p$-bit truth vector of a point. Let $m_M(c)$ count its occurrences in a nonempty finite structure $M$. For $q\ge0$ put $h^M_q(c)=\min(q,m_M(c))$.

\begin{theorem}[A complete monadic invariant]\label{thm:monadic}
Two such finite structures agree on every sentence of quantifier rank at most $q$ exactly when their nullary bits and all their $h_q(c)$ agree. From a failed comparison one can produce a distinguishing sentence of rank at most $q$. Finite satisfiability and finite validity in this language are PR decidable. Every finitely satisfiable sentence of rank $q\ge1$ has a model of at most $q2^p$ points, a bound attained for every $q\ge1,p\ge0$. At rank zero a one-point model suffices.
\end{theorem}
\begin{proof}
In the $q$-round game, match previously chosen points to their partners. For a new point choose an unused point of the same colour. With $r$ rounds remaining maintain equality of the tuples' equality patterns and colours, and, for each colour, equality of unused counts truncated at $r$. If the chosen colour has a new point on one side, the other has one because $r>0$. When the counts below $r$ coincide they both decrease by one, and when both are at least $r$ the new counts are both at least $r-1$. Repeating a point only lowers the threshold. At the leaves all relation atoms and equalities on the tuple agree. This is the game of Section~\ref{sec:ef}. Bounded search selects the least eligible reply. Alternatively the same argument transfers an existential witness and a universal counterexample at each quantifier.

For the converse, a mismatching nullary atom distinguishes immediately. If the smaller truncated count is $k$ and the larger one exceeds $k$, let $D_c(x)$ be the conjunction prescribing all $p$ bits, using $x=x$ for an empty conjunction. Then
\[
 \exists x_0\cdots\exists x_k\,
  \Big(\bigwedge_{i\le k}D_c(x_i)\ \wedge\
                  \bigwedge_{i<j\le k}x_i\ne x_j\Big)
\]
has rank $k+1\le q$ and holds exactly on the larger side. This follows by sorting or scanning for distinct elements of that colour. If an orientation is required, negate it. Hence equality of all bounded-rank truths, although expressed with an unbounded formula-code quantifier, has a PR complete test with a PR witness to failure.

For $q>0$ retain the first $h_q(c)$ elements of every colour, then reindex. The retained structure is nonempty and has the same invariant. To decide, enumerate the $2^z(q+1)^{2^p}$ possible bit/count profiles and discard the all-zero count vector. Realize each by its listed number of points of each colour and evaluate the sentence. For $q=0$ take each nullary valuation and one point of arbitrary colour. To see sharpness, conjoin for all $c$ the assertion that at least $q$ distinct points have colour $c$, using a separate block of $q$ existential binders inside each conjunct. Its rank is $q$. Unequal colours cannot share a point. A model with $q$ points of each colour realizes the sentence at the bound. In the equality-free language, matching a colour permits unlimited reuse of the same partner. Already one point per realized colour suffices at every positive rank.
\end{proof}

\paragraph{A normal form for monadic sentences.}
For fixed $q>0$ the conjunction describing a profile specifies the nullaries and, for each colour of count $k<q$, asserts at least $k$ and fewer than $k+1$ points. For count $q$ it asserts at least $q$. The zero threshold means truth. All parts have rank at most $q$. Disjoin these conjunctions over the accepting profiles. By the theorem this Boolean combination of threshold assertions is equivalent to the original sentence on every finite input. The empty disjunction means falsity. At rank zero use the truth table of nullaries alone. Classically the same decision answers unrestricted monadic satisfiability, because an arbitrary structure admits the same truncation of finitely many colour classes. For an infinite presented domain, deciding whether an individual colour occurs need not be effective. The internal uniform algorithm acts on sentences and finite tables.

\begin{theorem}[Proof-producing monadic decision]\label{thm:monadiccert}
For a closed monadic sentence $A$ in the preceding signature, a PR procedure returns either a finite model of $A$ or a classical deduction of $\neg A$. Applying it to $\neg A$ decides classical theoremhood in this fragment with a proof or a finite countermodel.
\end{theorem}
\begin{proof}
For each colour $c$, write $H_{c,j}$ for the closed threshold formula just used (with $j$ witnesses, and $H_{c,0}=\top$). We describe quantifier elimination to a Boolean combination of nullaries, tests on free variables, and the $H_{c,j}$ treated as closed units. In addition to the term for the result the transformation constructs a classical proof of its equivalence to the input. Atom and Boolean steps use the corresponding logical rules. Replace a universal by $\neg\exists\neg$, with its classical proof.

It remains to eliminate $\exists x\,D$ after the child has been translated. Expand $D$ into disjunctive normal form treating a threshold as one propositional atom. The propositional equivalence has a proof by splitting on its finitely many atoms. At a full valuation conjunction, disjunction, implication and negation are introduced or refuted from their signed children. Distribute existence over the disjunction and work on one conjunction. First orient equality and inequality involving $x$ by symmetry. Remove occurrences of $x=x$ and discard any conjunctive case containing $x\ne x$. If some conjunct requires $x=y_i$, substitute $y_i$ throughout, remove the quantifier, and retain the resulting conjunction. In the forward direction eliminate the witness and transport along the required equality. In the reverse direction use $y_i$ as witness. If no such equality occurs, all the remaining conditions on $x$ prescribe or prohibit unary bits, and require inequality from a list of parameters.

Split on every colour of the remaining parameters, and every equality between them, using excluded middle on the corresponding atoms. Inconsistent patterns give contradictions by equality transport. Split also on the possible colour $c$ of the witness. These alternatives are exhaustive by the $p$ atomic cases. Discard colours inconsistent with the unary conditions on $x$. In a surviving case let $d$ be the number of distinct forbidden parameters of colour $c$, calculated from the equality pattern. Replace the existence of $x$ by $H_{c,d+1}$, retaining the case conditions and all literals independent of $x$. Indeed a witness and one representative of each forbidden equality class supply $d+1$ pairwise distinct points of colour $c$. Conversely, open $d+1$ distinct witnesses for the threshold. Compare them to the $d$ forbidden representatives. If all were forbidden, their matching indices would be $d+1$ distinct numbers below $d$, impossible by the elementary marking argument. The comparisons give a surviving witness by a finite decision tree of equality cases. Its colour already ensures inequality from the forbidden parameters of other colours. This gives an actual classical proof template at each finite $d$, generated by expanding those cases, including $d=0$.

For a closed input, when eliminating a binder under $h$ other binders its remaining free parameters number at most $h$, so every new threshold is at most $h+1\le q$. Closed threshold units introduced earlier retain their scope. At the root, decide all nullaries and all thresholds up to $q$ by classical case splits inside the desired proof. Derive $H_{c,j+1}\to H_{c,j}$ by omitting the last witness, which removes all nonmonotone bit patterns. When $q>0$ and all $H_{c,1}$ are denied, take a fresh individual, determine its colour and introduce it as a witness to contradict the corresponding denial. Parameter erasure closes that branch. All remaining profiles have the canonical models of Theorem~\ref{thm:monadic}. At $q=0$ use one point. If some model accepts, output it. Otherwise the root Boolean formula is false in each consistent profile, and its signed atomic assumptions prove its negation by propositional evaluation. Recombine the cases and use the compiled equivalence to conclude $\neg A$. Running on $\neg A$ and eliminating double negation gives the theoremhood assertion by finite soundness.

All DNF, pattern and witness-comparison stages enumerate finite bitstrings with lengths computed from the current expression, and each stage constructs a finite rule tree. Treat introduced thresholds as marked units while processing parents, expanding them only in the output proofs. A fixed PR majorant for one such stage can be iterated once per original syntax node. Structural correctness collects the local proof-check bits as in Lemma~\ref{lem:toolkit}.
\end{proof}

\subsection{Two finite graphs a sentence cannot distinguish}
Connectedness and reachability are explicit examples in the starting article. Both already fail to be definable over finite undirected graphs. A direct construction here avoids any appeal to compactness.

\begin{theorem}[Rank-bounded witnesses against reachability]\label{thm:cycles}
Given $q$, put $s=2^q$ and $L=8(q+3)(s+1)$. Let $G$ be a cycle of length $2L$, with distinguished vertices $0,L$. Let $H$ be two disjoint cycles of length $L$, with one distinguished vertex in each. These pointed graphs agree at the displayed pair on every formula of rank at most $q$ whose free variables lie in those two slots, in the language of adjacency and equality. The distinguished vertices are joined by a path in $G$ and lie in different components of $H$. Even without distinguished vertices $G$ and $H$ agree through rank $q$.
\end{theorem}
\begin{proof}
Orient each cycle in its cyclic coordinate system, solely to describe Duplicator's replies. At $r$ rounds remaining preserve, for every pair of marked vertices, every displacement $d$ with $|d|\le 2^r$: the second is the first plus $d$ in its cycle on one side exactly when this holds on the other. Cross-component displacements are always false. These intervals give unique displacements whenever one exists, by the lower bound on the cycle lengths. The two initial marks satisfy the invariant including displacement zero.

Put $t=2^{r-1}$ for a move with $r>0$. If the new vertex is an offset $d$, $|d|\le t$, from an old mark, return the same offset from the corresponding mark. Any offset $e$ of length at most $t$ between the new vertex and another mark gives an offset of length at most $2t$ between the two old marks. Its truth transfers in both directions and subtraction of $d$ gives exactly the desired offset. If the vertex is far from all old marks, choose on the other side the least vertex outside all their radius-$t$ intervals. There are at most $q+2$ old marks, each excluding at most $2t+1$ vertices. Their total is bounded by $(q+2)(2t+1)<2L$. A bounded scan therefore succeeds. All old incidences retain the narrower invariant, and the new vertex has just its zero self-offset in the range. The argument works whichever board was chosen. At $r=0$ the displacements $0,1,-1$ give equality and undirected adjacency. Rank invariance now gives the assertion. The same strategy ignoring an unused pair of initial marks gives the sentence version. Traversing $L$ positive steps in $G$ gives the path. In $H$ an induction on path length preserves the component label. A proposed defining formula codes its own rank and hence its own counterexample pair.
\end{proof}
